\textbf{Phần 1. Trắc nghiệm lựa chọn}
\setcounterpageref{ex}{0}

\begin{ex} %[4D8V43]
Kết quả của phép nhân $34 \times 12$ là
\choice
{$406$}
{$418$}
{$398$}
{\True $408$}
\loigiai{
$34 \times 12 = 408$.\\
\textbf{Đáp án của câu hỏi là D}
}
\end{ex}

\begin{ex} %[4D8V43]
Trong phép nhân với số có hai chữ số, tích riêng thứ hai được viết
\choice
{Thẳng cột với tích riêng thứ nhất}
{\True Lùi sang bên trái một cột so với tích riêng thứ nhất}
{Lùi sang bên phải một cột}
{Lùi sang bên trái hai cột}
\loigiai{
Tích riêng thứ hai lùi sang bên trái một cột.\\
\textbf{Đáp án của câu hỏi là B}
}
\end{ex}

\begin{ex} %[4D8V43]
Kết quả của phép nhân $125 \times 24$ là
\choice
{$2800$}
{$3200$}
{\True $3000$}
{$3500$}
\loigiai{
$125 \times 24 = 3000$.\\
\textbf{Đáp án của câu hỏi là C}
}
\end{ex}

\begin{ex} %[4D8V43]
Kết quả của phép tính $245 \times 36$ là
\choice
{\True $8820$}
{$8720$}
{$8920$}
{$8620$}
\loigiai{
$245 \times 36 = 8820$.\\
\textbf{Đáp án của câu hỏi là A}
}
\end{ex}

\begin{ex} %[4D8V43]
Một quyển sách dày 148 trang. Hỏi 25 quyển sách như thế có tất cả bao nhiêu trang?
\choice
{$3600$ trang}
{$3800$ trang}
{$3500$ trang}
{\True $3700$ trang}
\loigiai{
$148 \times 25 = 3700$ trang.\\
\textbf{Đáp án của câu hỏi là D}
}
\end{ex}

\begin{ex} %[4D8V43]
Một hình chữ nhật có chiều dài $45\text{ m}$, chiều rộng $28\text{ m}$. Diện tích hình chữ nhật là
\choice
{$1240\text{ m}^2$}
{\True $1260\text{ m}^2$}
{$1280\text{ m}^2$}
{$1300\text{ m}^2$}
\loigiai{
$45 \times 28 = 1260\text{ m}^2$.\\
\textbf{Đáp án của câu hỏi là B}
}
\end{ex}

\begin{ex} %[4D8V43]
Một người thợ làm việc trong 24 ngày, mỗi ngày làm được 15 sản phẩm. Tổng số sản phẩm người đó làm được là
\choice
{$340$ sản phẩm}
{$380$ sản phẩm}
{\True $360$ sản phẩm}
{$400$ sản phẩm}
\loigiai{
$24 \times 15 = 360$ sản phẩm.\\
\textbf{Đáp án của câu hỏi là C}
}
\end{ex}

\begin{ex} %[4D8V43]
Một đoàn tàu hỏa có 12 toa xe, mỗi toa xe chở $35$ tấn hàng. Tổng khối lượng hàng đoàn tàu chở là
\choice
{\True $420\text{ tấn}$}
{$400\text{ tấn}$}
{$440\text{ tấn}$}
{$450\text{ tấn}$}
\loigiai{
$12 \times 35 = 420$ tấn.\\
\textbf{Đáp án của câu hỏi là A}
}
\end{ex}

\begin{ex} %[4D8V43]
Tính giá trị biểu thức: $245 \times 18 + 120 \times 25$.
\choice
{$7310$}
{$7510$}
{$7210$}
{\True $7410$}
\loigiai{
$4410 + 3000 = 7410$.\\
\textbf{Đáp án của câu hỏi là D}
}
\end{ex}

\begin{ex} %[4D8V43]
Một cửa hàng bán được 35 chiếc xe đạp với giá $1\,250\,000$ đồng một chiếc. Tổng số tiền bán xe đạp là
\choice
{$42\,750\,000$ đồng}
{\True $43\,750\,000$ đồng}
{$44\,750\,000$ đồng}
{$45\,000\,000$ đồng}
\loigiai{
$1\,250\,000 \times 35 = 43\,750\,000$ đồng.\\
\textbf{Đáp án của câu hỏi là B}
}
\end{ex}

\begin{ex} %[4D8V43]
Tìm $x$ biết: $x : 24 = 135$.
\choice
{$3140$}
{$3340$}
{\True $3240$}
{$3440$}
\loigiai{
$x = 135 \times 24 = 3240$.\\
\textbf{Đáp án của câu hỏi là C}
}
\end{ex}

\begin{ex} %[4D8V43]
Khi nhân một số tự nhiên với 35, một bạn học sinh đã đặt các tích riêng thẳng cột như trong phép cộng nên được kết quả là $680$. Tích đúng của phép nhân đó là
\choice
{\True $2975$}
{$2875$}
{$3075$}
{$3175$}
\loigiai{
Đặt tích riêng thẳng cột tức là nhân số đó với $3 + 5 = 8$.\\
Số đó là: $680 : 8 = 85$.\\
Tích đúng: $85 \times 35 = 2975$.\\
\textbf{Đáp án của câu hỏi là A}
}
\end{ex}

\textbf{Phần 2. Trắc nghiệm đúng sai}
\setcounterpageref{ex}{0}

\begin{ex} %[4D8V43]
Xét tính đúng sai của các phép tính:
\choiceTF
{\True $25 \times 14 = 350$.}
{$42 \times 23 = 946$.}
{\True $42 \times 23 = 966$.}
{\True Tích riêng thứ hai lùi sang trái một cột so với tích riêng thứ nhất.}
\loigiai{
\begin{itemize}
    \item $25 \times 14 = 350$. Đúng.
    \item $42 \times 23 = 966 \neq 946$. Ý 2 Sai, Ý 3 Đúng.
    \item Quy tắc đặt tính đúng.
\end{itemize}
\begin{itemize}
    \item a) \textbf{Đúng}.
    \item b) \textbf{Sai}.
    \item c) \textbf{Đúng}.
    \item d) \textbf{Đúng}.
\end{itemize}
}
\end{ex}

\begin{ex} %[4D8V43]
Một xí nghiệp may có 25 công nhân, mỗi ngày mỗi công nhân may được 12 chiếc áo. Mỗi chiếc áo bán với giá $150\,000$ đồng. Xét tính đúng sai:
\choiceTF
{\True Mỗi ngày xí nghiệp may được 300 chiếc áo.}
{\True Trong 10 ngày xí nghiệp may được 3000 chiếc áo.}
{Doanh thu mỗi ngày của xí nghiệp là 40 triệu đồng.}
{\True Doanh thu mỗi ngày của xí nghiệp là 45 triệu đồng.}
\loigiai{
\begin{itemize}
    \item $25 \times 12 = 300$ áo: Đúng.
    \item 10 ngày: $300 \times 10 = 3000$ áo: Đúng.
    \item Doanh thu: $300 \times 150\,000 = 45\,000\,000$ đồng $\neq 40$ triệu. Ý 3 Sai, Ý 4 Đúng.
\end{itemize}
\begin{itemize}
    \item a) \textbf{Đúng}.
    \item b) \textbf{Đúng}.
    \item c) \textbf{Sai}.
    \item d) \textbf{Đúng}.
\end{itemize}
}
\end{ex}

\textbf{Phần 3. Trả lời ngắn}
\setcounterpageref{ex}{0}

\begin{ex} %[4D8V43]
Tính tích: $36 \times 25$.
\shortans{900}
\loigiai{
\textbf{Đáp số:} $900$\\
$900$.
}
\end{ex}

\begin{ex} %[4D8V43]
Tính tích: $124 \times 15$.
\shortans{1860}
\loigiai{
\textbf{Đáp số:} $1860$\\
$1860$.
}
\end{ex}

\begin{ex} %[4D8V43]
Một mảnh vườn hình chữ nhật dài $35\text{ m}$, rộng $18\text{ m}$. Tính diện tích mảnh vườn theo mét vuông.
\shortans{630}
\loigiai{
\textbf{Đáp số:} $630$\\
$35 \times 18 = 630\text{ m}^2$.
}
\end{ex}

\begin{ex} %[4D8V43]
Một thùng sách có 45 cuốn, mỗi cuốn giá $24\,000$ đồng. Tính tổng giá trị thùng sách theo nghìn đồng (chỉ ghi số nghìn đồng).
\shortans{1080}
\loigiai{
\textbf{Đáp số:} $1080$\\
$45 \times 24 = 1080$ nghìn đồng.
}
\end{ex}

\begin{ex} %[4D8V43]
Tìm $x$ biết: $x : 16 = 250$.
\shortans{4000}
\loigiai{
\textbf{Đáp số:} $4000$\\
$x = 250 \times 16 = 4000$.
}
\end{ex}

\begin{ex} %[4D8V43]
Khi nhân một số với 24, bạn Bình quên lùi tích riêng thứ hai nên kết quả bị giảm đi $432$ đơn vị. Tìm số đem nhân với 24.
\shortans{24}
\loigiai{
\textbf{Đáp số:} $24$\\
Quên lùi tích riêng tức là đã nhân với $2 + 4 = 6$.\\
Tích giảm đi: $24 - 6 = 18$ lần số đó.\\
Số đó là: $432 : 18 = 24$.
}
\end{ex}
